\documentclass[12pt]{article}
\usepackage[utf8]{inputenc}
\usepackage[T1]{fontenc}
\usepackage{hyperref}
\title{Silopi Language Resources for K-12 Curriculum Development}
\author{Compiled by Claude Sonnet 5.5 (claude-sonnet-5-5) for Andr\'as Kornai}
\date{}
\begin{document}
\maketitle

\section{Introduction}
Silopi (ISO 639-3 xsp; Glottocode silo1240) is a Trans-New Guinea (Madang) language of Papua New Guinea. Speaker counts differ by source: Wikipedia and ASJP give 180, Global Recordings Network gives 140 and Joshua Project gives 300. Silopi should not be confused with Siliput (mkc). Very little is available and no teaching materials were found.

\section{Primary Educational Resources}
No primers, textbooks or school programmes were found.

\section{Linguistic and Reference Materials}
\begin{itemize}
\item Wikipedia, ``Silopi language''. \url{https://en.wikipedia.org/wiki/Silopi_language}
\item ASJP word list for Silopi, compiled by Cecil H. Brown from Z'graggen's 1980 Mabuso wordlist. \url{https://asjp.clld.org/languages/SILOPI}
\item Glottolog languoid page. The URL is the one returned by the search and carries a different id from the Glottocode given in the task, so check it before citing. \url{https://glottolog.org/resource/languoid/id/sili1247}
\item Endangered Languages Project bibliography. \url{https://www.endangeredlanguages.com/lang/3228/bibliography?hl=ru}
\end{itemize}

\section{Religious and Community Resources}
\begin{itemize}
\item Joshua Project, Silopi in Papua New Guinea. It reports no Bible translation. \url{https://joshuaproject.net/people_groups/14177/PP}
\item Joshua Project, Silopi language page. \url{https://joshuaproject.net/languages/xsp}
\item Scripture Earth page for xsp. I did not check what it holds. \url{https://www.scriptureearth.org/00eng.php?iso=xsp}
\end{itemize}

\section{Digital and Audio Resources}
Global Recordings Network has a page for xsp, but the search summary indicates no audio is currently available. \url{https://globalrecordings.net/en/language/xsp}

\section{Curriculum Development Context}
No evidence was found of Silopi in schools. The language is endangered with a few hundred speakers at most. Curriculum work would need new documentation with the community, for example orthography design, oral-literature recording and a vernacular elementary scheme, in cooperation with PNG education authorities. None of this was found in the search results.
\end{document}